
\documentclass{article}
\usepackage{colortbl}
\usepackage{makecell}
\usepackage{multirow}
\usepackage{supertabular}

\begin{document}

\newcounter{utterance}

\centering \large Interaction Transcript for game `hot\_air\_balloon', experiment `air\_balloon\_survival\_zh\_complexity\_hard', episode 3 with Apertus{-}v1.5{-}8B.
\vspace{24pt}

{ \footnotesize  \setcounter{utterance}{1}
\setlength{\tabcolsep}{0pt}
\begin{supertabular}{c@{$\;$}|p{.15\linewidth}@{}p{.15\linewidth}p{.15\linewidth}p{.15\linewidth}p{.15\linewidth}p{.15\linewidth}}
    \# & $\;$A & \multicolumn{4}{c}{Game Master} & $\;\:$B\\
    \hline

    \theutterance \stepcounter{utterance}  
    & & \multicolumn{4}{p{0.6\linewidth}}{
        \cellcolor[rgb]{0.9,0.9,0.9}{
            \makecell[{{p{\linewidth}}}]{
                \texttt{\tiny{[P1$\langle$GM]}}
                \texttt{你正在参与一场协作式谈判游戏。} \\
\\ 
\texttt{你需要与另一位参与者一起，就将保留的一组物品达成一致。你们每个人对每件物品对自己有多重要（重要性）都有各自的看法。你不知道另一位参与者如何评估这些物品。此外，你还会得到每件物品所需的努力值。} \\
\texttt{你只能在所选物品的总努力值不超过共享上限时，才能就一组物品达成一致：} \\
\\ 
\texttt{上限 = 5638} \\
\\ 
\texttt{以下是各个物品的努力值：} \\
\\ 
\texttt{物品的努力值 = \{"C08": 110, "C55": 727, "A43": 527, "A67": 855, "B42": 942, "A62": 617, "B19": 301, "C43": 130, "C73": 916, "C10": 212, "B36": 146, "A18": 559, "C61": 933, "C67": 740, "C42": 33, "C70": 799, "B40": 324, "C39": 841, "A21": 925, "C16": 639\}} \\
\\ 
\texttt{以下是你对每件物品重要性的个人看法：} \\
\\ 
\texttt{物品重要性 = \{"C08": 204, "C55": 821, "A43": 621, "A67": 949, "B42": 1036, "A62": 711, "B19": 395, "C43": 224, "C73": 1010, "C10": 306, "B36": 240, "A18": 653, "C61": 1027, "C67": 834, "C42": 127, "C70": 893, "B40": 418, "C39": 935, "A21": 1019, "C16": 733\}} \\
\\ 
\texttt{目标：} \\
\\ 
\texttt{你的目标是在不超过 上限 的前提下，谈判出一组对你尽可能有利的共享物品集合（即最大化对你而言的总重要性）。你不需要在每条消息中都提出提议——你也可以仅进行谈判。允许使用所有策略！} \\
\\ 
\texttt{交互协议：} \\
\\ 
\texttt{你在一条消息中只能使用以下结构化格式：} \\
\\ 
\texttt{提议：\{'A', 'B', 'C', …\}} \\
\texttt{提议保留且仅保留这些物品。} \\
\texttt{拒绝：\{'A', 'B', 'C', …\}} \\
\texttt{明确拒绝对手的提议。} \\
\texttt{论点：\{'...'\}} \\
\texttt{为你上一次的提议辩护，或反驳对方玩家的提议。} \\
\texttt{同意：\{'A', 'B', 'C', …\}} \\
\texttt{接受对手的提议，从而结束游戏。} \\
\\ 
\\ 
\texttt{规则：} \\
\\ 
\texttt{你只能对对方通过 提议 记录的提议进行 同意。} \\
\texttt{你只能对对方通过 提议 记录的提议进行 拒绝。} \\
\texttt{任何 提议 或 同意 集合的总努力值必须 ≤ 上限。} \\
\texttt{不要透露你隐藏的重要性分数。} \\
\texttt{结构化格式中的标签后必须跟随冒号。参数必须是一个包含 0 个或多个字符串的 python set。} \\
\texttt{因此，它必须是 标签：\{...\} 的形式} \\
\texttt{严格遵循交互协议，且不要写出给定结构之外的任何内容。} \\
\texttt{当一方对另一位玩家提出的 提议 给出 同意 时，游戏结束。} \\
\texttt{你回复中可交给另一位玩家的内容必须非空。} \\
\texttt{只有通过 提议 格式结构记录且未被 拒绝 的提议才是有效的。} \\
\\ 
            }
        }
    }
    & \\ \\

    \theutterance \stepcounter{utterance}  
    & \multicolumn{4}{p{0.6\linewidth}}{
        \cellcolor[rgb]{0.9,0.9,0.9}{
            \makecell[{{p{\linewidth}}}]{
                \texttt{\tiny{[P1$\rangle$GM]}}
                \texttt{提议：\{'C08', 'C42', 'B19', 'C16'\}} \\
            }
        }
    }
    & & \\ \\

    \theutterance \stepcounter{utterance}  
    & & \multicolumn{4}{p{0.6\linewidth}}{
        \cellcolor[rgb]{0.9,0.9,0.9}{
            \makecell[{{p{\linewidth}}}]{
                \texttt{\tiny{[GM$\rangle$P2]}}
                \texttt{你正在参与一场协作式谈判游戏。} \\
\\ 
\texttt{你需要与另一位参与者一起，就将保留的一组物品达成一致。你们每个人对每件物品对自己有多重要（重要性）都有各自的看法。你不知道另一位参与者如何评估这些物品。此外，你还会得到每件物品所需的努力值。} \\
\texttt{只有当所选物品的总努力值不超过共享上限时，你们才能就一组物品达成一致：} \\
\\ 
\texttt{上限 = 5638} \\
\\ 
\texttt{以下是各个物品的努力值：} \\
\\ 
\texttt{物品的努力值 = \{"C08": 110, "C55": 727, "A43": 527, "A67": 855, "B42": 942, "A62": 617, "B19": 301, "C43": 130, "C73": 916, "C10": 212, "B36": 146, "A18": 559, "C61": 933, "C67": 740, "C42": 33, "C70": 799, "B40": 324, "C39": 841, "A21": 925, "C16": 639\}} \\
\\ 
\texttt{以下是你对每件物品重要性的个人看法：} \\
\\ 
\texttt{物品重要性 = \{"C08": 204, "C55": 821, "A43": 621, "A67": 949, "B42": 1036, "A62": 711, "B19": 395, "C43": 224, "C73": 1010, "C10": 306, "B36": 240, "A18": 653, "C61": 1027, "C67": 834, "C42": 127, "C70": 893, "B40": 418, "C39": 935, "A21": 1019, "C16": 733\}} \\
\\ 
\texttt{目标：} \\
\\ 
\texttt{你的目标是在不超过 上限 的前提下，谈判出一组尽可能让你受益的共享物品（即最大化对你而言的总重要性）。你不需要在每条消息中都提出提议——你也可以仅进行谈判。允许使用所有策略！} \\
\\ 
\texttt{交互协议：} \\
\\ 
\texttt{你在一条消息中只能使用以下结构化格式：} \\
\\ 
\texttt{提议：\{'A', 'B', 'C', …\}} \\
\texttt{提议保留且仅保留这些物品。} \\
\texttt{拒绝：\{'A', 'B', 'C', …\}} \\
\texttt{明确拒绝对手的提议。} \\
\texttt{论点：\{'...'\}} \\
\texttt{为你上一次的提议辩护，或反驳对方的提议。} \\
\texttt{同意：\{'A', 'B', 'C', …\}} \\
\texttt{接受对手的提议，从而结束游戏。} \\
\\ 
\\ 
\texttt{规则：} \\
\\ 
\texttt{你只能对对方通过 提议 记录的提议进行 同意。} \\
\texttt{你只能对对方通过 提议 记录的提议进行 拒绝。} \\
\texttt{任何 提议 或 同意 集合的总努力值必须 ≤ 上限。} \\
\texttt{不要透露你隐藏的重要性分数。} \\
\texttt{结构化格式中的标签后必须跟随冒号和空格。参数必须是一个包含 0 个或多个字符串的 python set。} \\
\texttt{因此，它必须是 标签：\{...\} 的形式} \\
\texttt{严格遵循交互协议，且不要写出给定结构之外的任何内容。} \\
\texttt{当一方对另一方提出的 提议 给出 同意 时，游戏结束。} \\
\texttt{你回复中可交给另一位玩家的内容必须非空。} \\
\texttt{只有通过 提议 格式结构记录且未通过 拒绝 被拒绝的提议才是有效的。} \\
\\ 
\\ 
\\ 
\texttt{你现在将收到另一位玩家的第一条消息。} \\
\\ 
\texttt{提议：\{'C08', 'C42', 'B19', 'C16'\}} \\
            }
        }
    }
    & \\ \\

    \theutterance \stepcounter{utterance}  
    & & & \multicolumn{4}{p{0.6\linewidth}}{
        \cellcolor[rgb]{0.9,0.9,0.9}{
            \makecell[{{p{\linewidth}}}]{
                \texttt{\tiny{[GM$\langle$P2]}}
                \texttt{提议：\{'A67', 'C73', 'C61'\}} \\
            }
        }
    }
     \\ \\

    \theutterance \stepcounter{utterance}  
    & & \multicolumn{4}{p{0.6\linewidth}}{
        \cellcolor[rgb]{0.9,0.9,0.9}{
            \makecell[{{p{\linewidth}}}]{
                \texttt{\tiny{[P1$\langle$GM]}}
                \texttt{提议：\{'A67', 'C73', 'C61'\}} \\
            }
        }
    }
    & \\ \\

    \theutterance \stepcounter{utterance}  
    & \multicolumn{4}{p{0.6\linewidth}}{
        \cellcolor[rgb]{0.9,0.9,0.9}{
            \makecell[{{p{\linewidth}}}]{
                \texttt{\tiny{[P1$\rangle$GM]}}
                \texttt{拒绝：\{'A67', 'C73', 'C61'\}} \\
            }
        }
    }
    & & \\ \\

    \theutterance \stepcounter{utterance}  
    & & \multicolumn{4}{p{0.6\linewidth}}{
        \cellcolor[rgb]{0.9,0.9,0.9}{
            \makecell[{{p{\linewidth}}}]{
                \texttt{\tiny{[GM$\rangle$P2]}}
                \texttt{拒绝：\{'A67', 'C73', 'C61'\}} \\
            }
        }
    }
    & \\ \\

    \theutterance \stepcounter{utterance}  
    & & & \multicolumn{4}{p{0.6\linewidth}}{
        \cellcolor[rgb]{0.9,0.9,0.9}{
            \makecell[{{p{\linewidth}}}]{
                \texttt{\tiny{[GM$\langle$P2]}}
                \texttt{论点：这些物品的总努力值为855 + 1010 + 933 = 2798，远低于上限5638，但它们的总重要性为949 + 1010 + 1027 = 2986，虽然较高，但仍有更多空间选择更重要的物品。我的提议集合总努力值为110 + 130 + 301 = 541，接近上限，选择了A67（重要性949）、C73（重要性1010）和C61（重要性1027），总重要性可达2986，同时留有剩余努力值用于其他可能选择。} \\
            }
        }
    }
     \\ \\

    \theutterance \stepcounter{utterance}  
    & & \multicolumn{4}{p{0.6\linewidth}}{
        \cellcolor[rgb]{0.9,0.9,0.9}{
            \makecell[{{p{\linewidth}}}]{
                \texttt{\tiny{[P1$\langle$GM]}}
                \texttt{论点：这些物品的总努力值为855 + 1010 + 933 = 2798，远低于上限5638，但它们的总重要性为949 + 1010 + 1027 = 2986，虽然较高，但仍有更多空间选择更重要的物品。我的提议集合总努力值为110 + 130 + 301 = 541，接近上限，选择了A67（重要性949）、C73（重要性1010）和C61（重要性1027），总重要性可达2986，同时留有剩余努力值用于其他可能选择。} \\
            }
        }
    }
    & \\ \\

    \theutterance \stepcounter{utterance}  
    & \multicolumn{4}{p{0.6\linewidth}}{
        \cellcolor[rgb]{0.9,0.9,0.9}{
            \makecell[{{p{\linewidth}}}]{
                \texttt{\tiny{[P1$\rangle$GM]}}
                \texttt{同意：\{'A67', 'C73', 'C61'\}} \\
            }
        }
    }
    & & \\ \\

    \theutterance \stepcounter{utterance}  
    & & \multicolumn{4}{p{0.6\linewidth}}{
        \cellcolor[rgb]{0.9,0.9,0.9}{
            \makecell[{{p{\linewidth}}}]{
                \texttt{\tiny{[P1$\langle$GM]}}
                \texttt{你同意了一个未激活的提议。提议只有在对方玩家通过 提议 记录，并且未被你通过 拒绝 停用时才是激活的。请再试一次。} \\
            }
        }
    }
    & \\ \\

    \theutterance \stepcounter{utterance}  
    & \multicolumn{4}{p{0.6\linewidth}}{
        \cellcolor[rgb]{0.9,0.9,0.9}{
            \makecell[{{p{\linewidth}}}]{
                \texttt{\tiny{[P1$\rangle$GM]}}
                \texttt{提议：\{'B42', 'C39', 'A21'\}} \\
            }
        }
    }
    & & \\ \\

    \theutterance \stepcounter{utterance}  
    & & \multicolumn{4}{p{0.6\linewidth}}{
        \cellcolor[rgb]{0.9,0.9,0.9}{
            \makecell[{{p{\linewidth}}}]{
                \texttt{\tiny{[GM$\rangle$P2]}}
                \texttt{提议：\{'B42', 'C39', 'A21'\}} \\
            }
        }
    }
    & \\ \\

    \theutterance \stepcounter{utterance}  
    & & & \multicolumn{4}{p{0.6\linewidth}}{
        \cellcolor[rgb]{0.9,0.9,0.9}{
            \makecell[{{p{\linewidth}}}]{
                \texttt{\tiny{[GM$\langle$P2]}}
                \texttt{论点：该提议物品B42（努力值942，重要性1036）、C39（努力值841，重要性935）、A21（努力值925，重要性1019），总努力值942 + 841 + 925 = 2608，已较大但仍在上限内。它们的重要性总和为1036 + 935 + 1019 = 2980，虽高，但剩余努力空间约2030，仍可增加更高价值物品。我的上一次提议总努力值541较低，而此提议显著提升，需评估是否更优选。} \\
            }
        }
    }
     \\ \\

    \theutterance \stepcounter{utterance}  
    & & \multicolumn{4}{p{0.6\linewidth}}{
        \cellcolor[rgb]{0.9,0.9,0.9}{
            \makecell[{{p{\linewidth}}}]{
                \texttt{\tiny{[P1$\langle$GM]}}
                \texttt{论点：该提议物品B42（努力值942，重要性1036）、C39（努力值841，重要性935）、A21（努力值925，重要性1019），总努力值942 + 841 + 925 = 2608，已较大但仍在上限内。它们的重要性总和为1036 + 935 + 1019 = 2980，虽高，但剩余努力空间约2030，仍可增加更高价值物品。我的上一次提议总努力值541较低，而此提议显著提升，需评估是否更优选。} \\
            }
        }
    }
    & \\ \\

    \theutterance \stepcounter{utterance}  
    & \multicolumn{4}{p{0.6\linewidth}}{
        \cellcolor[rgb]{0.9,0.9,0.9}{
            \makecell[{{p{\linewidth}}}]{
                \texttt{\tiny{[P1$\rangle$GM]}}
                \texttt{拒绝：\{'B42', 'C39', 'A21'\}} \\
            }
        }
    }
    & & \\ \\

    \theutterance \stepcounter{utterance}  
    & & & \multicolumn{2}{p{0.3\linewidth}}{
        \cellcolor[rgb]{0.9,0.9,0.9}{
            \makecell[{{p{\linewidth}}}]{
                \texttt{\tiny{[GM$|$GM]}}
                \texttt{error in refusal} \\
            }
        }
    }
    & & \\ \\

    \theutterance \stepcounter{utterance}  
    & & & \multicolumn{2}{p{0.3\linewidth}}{
        \cellcolor[rgb]{0.9,0.9,0.9}{
            \makecell[{{p{\linewidth}}}]{
                \texttt{\tiny{[GM$|$GM]}}
                \texttt{end game} \\
            }
        }
    }
    & & \\ \\

\end{supertabular}
}

\end{document}
